\documentclass[../main.tex]{subfiles}
\begin{document}

\chapter{Synchronization}
\lessoninfo{
  video={Lesson 20, videos 1--23},
  textbook={H\&P \cite{hennessy2017} §5.5, App.~I},
  keywords={lock, mutual exclusion, critical section, spin lock, Lamport's bakery algorithm, atomic instruction, test-and-set, load-linked/store-conditional, link register, cache coherence, barrier synchronization},
}

Lesson 19 made the caches of all cores present one coherent view of the
shared memory.  Coherent memory alone does not make a shared-memory program
correct: when several threads read and update the same variable, the result
depends on how their instructions interleave.  This lesson introduces locks,
which admit only one thread at a time to such an update, shows why a lock
cannot be built from ordinary loads and stores and which atomic instructions
processors provide instead, examines how the lock variable interacts with
cache coherence, and ends with barriers, which make a group of threads wait
for each other.

\section{The need for synchronization}
\label{sec:l20-need}

Threads that share memory must coordinate their updates of shared
data.\source{Lesson 20, videos 1--2; H\&P §5.5.}  As an example, consider a
program that counts how often each letter occurs in a long document.  Two
threads each process one half of the document and increment a shared array
of 26 counters (\cref{ex:l20-count}).

\begin{example}[Counting letters]
  \label{ex:l20-count}
  Each thread executes the following loop over its half of the document,
  which contains only lowercase letters:
  \begin{lstlisting}[language=C, numbers=none, morekeywords={shared}]
shared char doc[N];
shared int  count[26];
for (int i = lo; i < hi; i++) {  /* my half */
  int c = doc[i] - 'a';
  count[c]++;
}
\end{lstlisting}
\end{example}

The statement \texttt{count[c]++} is one line of source code but three
instructions for the processor: a load of the counter, an addition and a
store.
\begin{lstlisting}[language={[x86masm]Assembler}, numbers=none]
LW   R1, 0(R2)      ; R2 = address of count[c]
ADDI R1, R1, 1
SW   R1, 0(R2)
\end{lstlisting}
Each thread has its own registers, but when both threads encounter the same
letter, their R2 hold the same address.  If both threads load the counter
before either stores the incremented value, both compute the same new value
and one increment is lost (\cref{tab:l20-race}).

\begin{table}[htbp]
  \centering\small
  \caption{Both threads increment the counter of the letter \texttt{e},
    which holds 41.  After two increments it holds 42 instead of 43.}
  \label{tab:l20-race}
  \begin{tabular}{@{}cllrrr@{}}
    \toprule
    Step & Thread A & Thread B & R1 of A & R1 of B & Counter \\
    \midrule
    1 & \texttt{LW R1, 0(R2)}   &                             & 41 & --  & 41 \\
    2 &                         & \texttt{LW R1, 0(R2)}       & 41 & 41  & 41 \\
    3 & \texttt{ADDI R1, R1, 1} &                             & 42 & 41  & 41 \\
    4 &                         & \texttt{ADDI R1, R1, 1}     & 42 & 42  & 41 \\
    5 & \texttt{SW R1, 0(R2)}   &                             & 42 & 42  & 42 \\
    6 &                         & \texttt{SW R1, 0(R2)}       & 42 & 42  & 42 \\
    \bottomrule
  \end{tabular}
\end{table}

Whether an increment is lost depends on the relative timing of the threads.
Most of the time the two threads process different letters, or one thread
finishes its three instructions before the other starts, and the counts are
correct; occasionally a count comes out too low, and the error need not
reappear when the program is run again.\margin{The same loss occurs on a
single time-shared core if a context switch falls between the load and the
store.}  The load, the addition and the store must therefore execute as one
indivisible unit with respect to every other thread that updates the same
counter.  Providing this is the task of \emph{synchronization}.

\section{Locks}
\label{sec:l20-locks}

The basic synchronization mechanism is the lock.\source{Lesson 20, videos
3--5; H\&P §5.5.}

\begin{definition}[Lock]
  A \term{lock}[ロック] is a shared object with two operations.
  \texttt{lock} waits until the lock is free and then marks it as held by the
  calling thread; \texttt{unlock} marks it free again.  At any time at most
  one thread holds a given lock.
\end{definition}

With one lock per counter, the loop of \cref{ex:l20-count} becomes
\begin{lstlisting}[language=C, numbers=none, morekeywords={shared}]
shared lock_t countLock[26];
  ...
  lock(&countLock[c]);
  count[c]++;
  unlock(&countLock[c]);
\end{lstlisting}
If thread B calls \texttt{lock} while thread A holds the lock, B waits until
A calls \texttt{unlock}, and only then loads the counter
(\cref{fig:l20-lock-timing}).  Thread B therefore sees the value that
thread~A has stored, and no increment is lost.

\begin{figure}[htbp]
  \centering
  % Two threads that use the same lock.  Thread B calls lock() while thread A
% is in its critical section, waits, and enters when A calls unlock().
\begin{tikzpicture}[x=0.9cm,y=1cm, font=\footnotesize\sffamily, >={Stealth[length=1.6mm]},
  ev/.style={font=\scriptsize\ttfamily, inner sep=1pt}]
  \def\lnbox#1#2#3#4{% y, x0, x1, fill
    \draw[fill=#4] (#2,{#1-0.25}) rectangle (#3,{#1+0.25});}
  % ---- thread A
  \node[anchor=east] at (-0.15,1) {\iflangja{スレッド A}{thread A}};
  \lnbox{1}{0}{1.5}{white}
  \lnbox{1}{1.5}{4.5}{ln-accent!20}
  \lnbox{1}{4.5}{10}{white}
  \node[font=\scriptsize] at (3,1) {\iflangja{クリティカルセクション}{critical section}};
  \node[ev, anchor=south] at (1.5,1.3) {lock()};
  \node[ev, anchor=south] at (4.5,1.3) {unlock()};
  \draw (1.5,1.25) -- (1.5,1.32);
  \draw (4.5,1.25) -- (4.5,1.32);
  % ---- thread B
  \node[anchor=east] at (-0.15,0) {\iflangja{スレッド B}{thread B}};
  \lnbox{0}{0}{2.5}{white}
  \lnbox{0}{2.5}{4.5}{ln-caution!15}
  \lnbox{0}{4.5}{7.5}{ln-accent!20}
  \lnbox{0}{7.5}{10}{white}
  \node[font=\scriptsize] at (3.5,0) {\iflangja{待機}{waiting}};
  \node[font=\scriptsize] at (6,0) {\iflangja{クリティカルセクション}{critical section}};
  \node[ev, anchor=north] at (2.5,-0.3) {lock()};
  \node[ev, anchor=north] at (7.5,-0.3) {unlock()};
  \draw (2.5,-0.25) -- (2.5,-0.32);
  \draw (7.5,-0.25) -- (7.5,-0.32);
  % ---- hand-over of the lock
  \draw[->, ln-caution] (4.5,0.75) -- (4.5,0.25);
  % ---- time axis
  \draw[->, ln-muted] (0,-0.95) -- (10.3,-0.95)
    node[right, inner sep=2pt] {\iflangja{時間}{time}};
\end{tikzpicture}

  \caption{Two threads using the same lock: thread B waits in
    \texttt{lock} until thread A leaves its critical section.}
  \label{fig:l20-lock-timing}
\end{figure}

\begin{definition}[Mutual exclusion]
  The guarantee that at most one thread at a time executes code protected
  by a given lock is called \term{mutual exclusion}[相互排除].
\end{definition}

For this reason a lock is also called a \emph{mutex}\index{mutex|see{lock}}.

\begin{definition}[Critical section]
  A \term{critical section}[クリティカルセクション] is a sequence of
  instructions that must execute under mutual exclusion; it begins with
  \texttt{lock} and ends with \texttt{unlock}.
\end{definition}

A lock protects \emph{data}, not code.  Mutual exclusion holds only among
threads that use the same lock, so every access to a shared variable must be
protected by one and the same lock.  How many variables one lock protects is
a choice with performance consequences:
\begin{itemize}
  \item \textbf{One lock per counter.}  Threads that increment different
    counters proceed in parallel; a thread waits only when another thread is
    updating the same counter.
  \item \textbf{One lock for all counters.}  This is also correct, but every
    increment excludes all others, even of different letters, and the
    threads spend much of their time waiting for each other.
  \item \textbf{A lock that depends on the thread}, for example a private
    lock per thread, is wrong: two threads updating the same counter hold
    different locks and do not exclude each other.
\end{itemize}
\tip{Use the same lock for every access to a given shared variable and
different locks for independent variables.}

\section{Implementing a lock}
\label{sec:l20-implementing}

A lock can be represented by a shared \emph{lock variable} that holds 0 when
the lock is free and 1 when it is held.\source{Lesson 20, videos 6--7; H\&P
§5.5; \cite{lamport1974}.}  The obvious implementation waits until the
variable is 0 and then sets it to~1:
\begin{lstlisting}[language=C, numbers=none]
void lock(lock_t *l) {
  while (*l == 1)
    ;              /* wait until the lock is free */
  *l = 1;          /* take it                     */
}
void unlock(lock_t *l) {
  *l = 0;
}
\end{lstlisting}
A waiting thread keeps executing the loop, reading the lock variable over
and over; a lock that waits in this way is a \term{spin lock}[スピンロック].

\needspace{7\baselineskip}
This implementation does not work.  In machine code, testing the lock
variable and setting it are separate instructions:
\begin{lstlisting}[language={[x86masm]Assembler}, numbers=none, deletekeywords={lock}]
lock:  LW   R1, 0(R2)    ; R2 = address of lock variable
       BNEZ R1, lock     ; held: read it again
       ADDI R1, R0, 1
       SW   R1, 0(R2)    ; lock variable <- 1
\end{lstlisting}
If two threads both execute the \texttt{LW} while the lock is free, both read
0, both leave the loop, both store 1, and both enter the critical section.
This is the lost update of \cref{sec:l20-need} once more: the test and the
set are a read and an update of a shared variable, so \texttt{lock} would
itself need a critical section---the very thing it is supposed to provide.

There are two ways out.
\begin{itemize}
  \item \textbf{Software.}  \term{Lamport's bakery algorithm}[Lamport のパン屋のアルゴリズム]
    \cite{lamport1974} achieves mutual exclusion with
    ordinary loads and stores.  A thread that wants to enter first takes a
    number larger than every number currently held by other threads, like a
    customer taking a ticket, and enters only when no other thread that wants
    to enter holds a smaller number (ties are broken by thread identifier).
    The algorithm is correct, but every acquisition must examine the
    variables of all threads, which takes many instructions---too slow for
    locks that are acquired very frequently.
  \item \textbf{Hardware.}  The processor provides an instruction that tests
    and sets the lock variable in one step.
\end{itemize}

\begin{definition}[Atomic instruction]
  An \term{atomic instruction}[アトミック命令] performs several steps---here,
  reading a memory location and writing it---as one indivisible operation:
  no other instruction, on this or any other core, can access the location
  between the steps.
\end{definition}

To implement a lock, the atomic step must contain both the read and the
write of the lock variable.  The test of the value that was read can follow
in ordinary instructions, because once the variable has been written, no
other thread can still read the old value.

\section{Test-and-set}
\label{sec:l20-tst}

The simplest atomic instruction for locking combines a load with a store of
a constant.\source{Lesson 20, videos 8--10; H\&P §5.5.}

\begin{definition}[Test-and-set]
  A \term{test-and-set}[テストアンドセット] instruction
  \texttt{TST R1, 0(R2)} atomically loads the memory location at the
  address in R2 into R1 and stores 1 into that location:
  $\text{R1} \leftarrow \text{Mem}[\text{R2}]$ and
  $\text{Mem}[\text{R2}] \leftarrow 1$.
\end{definition}

A thread acquires the lock by executing \texttt{TST} on the lock variable
until it returns 0.  If the lock is held, \texttt{TST} overwrites 1 with 1,
which changes nothing, and returns 1, so the thread tries again.  If the lock
is free, the first \texttt{TST} to execute returns 0 and leaves 1 behind, and
every later \texttt{TST} returns 1.  Because nothing can intervene between
the read and the write, two threads can never both read~0.\margin{H\&P also
describe an atomic \emph{exchange} of a register and a memory location; with
1 in the register it acts exactly like test-and-set.}

\needspace{6\baselineskip}
With test-and-set, \texttt{lock} and \texttt{unlock} thus take one
instruction each, plus a branch:
\begin{lstlisting}[language={[x86masm]Assembler}, numbers=none, deletekeywords={lock}]
lock:   TST  R1, 0(R2)   ; R1 <- lock var., lock var. <- 1
        BNEZ R1, lock    ; was held: try again
        ...              ; critical section
unlock: SW   R0, 0(R2)   ; lock variable <- 0
\end{lstlisting}

The drawback of test-and-set lies in its implementation.  In a load/store
instruction set every other instruction accesses data memory at most once:
a load reads and a store writes, both in the single MEM stage of the
pipeline (Lesson 3).  Test-and-set must read \emph{and} write the same
location within one instruction.  The pipeline needs a special datapath with
two memory accesses for this single instruction, and the processor must
ensure that nothing---neither an access by another core nor an
exception---comes between the two.

\section{Load-linked/store-conditional}
\label{sec:l20-llsc}

The alternative splits the atomic read and write into two instructions,
each of which accesses memory only once, like an ordinary load or
store.\source{Lesson 20, videos 11--13; H\&P §5.5.}

\begin{definition}[Load-linked/store-conditional]
  The instruction pair \term{load-linked/store-conditional}[LL/SC]%
  \index{LL/SC|see{load-linked/store-conditional}} (LL/SC) consists of two
  instructions:
  \begin{description}
    \item[\texttt{LL R1, 0(R2)}] loads like \texttt{LW},
      $\text{R1} \leftarrow \text{Mem}[\text{R2}]$, and records the address
      in the \term{link register}[リンクレジスタ] of the core;
    \item[\texttt{SC R1, 0(R2)}] stores like \texttt{SW},
      $\text{Mem}[\text{R2}] \leftarrow \text{R1}$, but only if the link
      register still holds this address; it then sets R1 to 1 to report
      success.  Otherwise it stores nothing and sets R1 to~0.
  \end{description}
\end{definition}

What makes the pair atomic is the rule for clearing the link register: every
write to a memory location---an \texttt{SC} or an ordinary store, by this
core or by any other---clears every link register that holds the address of
that location (\cref{fig:l20-llsc}).\margin{Per H\&P, an interrupt or
exception between \texttt{LL} and \texttt{SC} also clears the link
register.}  An \texttt{SC} therefore succeeds only if nobody has written the
location since the matching \texttt{LL}, that is, only if the value that
\texttt{LL} read is still the value in memory at the moment of the store.
Other instructions may run between \texttt{LL} and \texttt{SC}, including
instructions of other threads that access the same location; if they
interfere, the \texttt{SC} detects it and fails, and
the program simply tries again.  The effect is that of an atomic read and
write, although neither instruction by itself does more than a normal load
or store.

\needspace{9\baselineskip}
A lock built with LL/SC loads the lock variable, tests it, and tries to
store 1 only if it was~0:\tip{Keep the instructions between \texttt{LL}
and \texttt{SC} few and free of memory accesses: every extra instruction
widens the window in which a write by another thread makes the \texttt{SC}
fail.}
\begin{lstlisting}[language={[x86masm]Assembler}, numbers=none, deletekeywords={lock}]
lock:   LL   R1, 0(R2)   ; R1 <- lock variable, link
        BNEZ R1, lock    ; held: try again
        ADDI R1, R0, 1   ; R1 <- 1
        SC   R1, 0(R2)   ; store 1 if still linked
        BEQZ R1, lock    ; SC failed: start over
        ...              ; critical section
unlock: SW   R0, 0(R2)   ; lock variable <- 0
\end{lstlisting}
\begin{figure}[htbp]
  \centering
  % LL/SC on two cores.  Both cores load-link the lock variable at address A;
% core 0's store-conditional succeeds, and its store clears every link
% register that holds A, so core 1's store-conditional fails.
% Register contents are shown after step (5).
\begin{tikzpicture}[x=1cm,y=1cm, font=\footnotesize\sffamily, >={Stealth[length=1.8mm]},
  core/.style={draw, fill=ln-accent!15, minimum width=3.6cm, minimum height=0.55cm, inner sep=2pt},
  reg/.style={draw, fill=white, minimum width=1.8cm, minimum height=0.5cm, inner sep=1pt, font=\footnotesize\ttfamily},
  hd/.style={inner sep=1pt, font=\scriptsize\sffamily, text=ln-muted},
  big/.style={draw, fill=ln-box, minimum height=0.65cm, inner sep=3pt},
  op/.style={font=\scriptsize, align=left, anchor=south west, inner sep=1pt}]
  % ---- core 0
  \node[core] (c0) at (1.8,0) {\iflangja{コア}{core}~0};
  \node[reg, anchor=north west] (r0) at (c0.south west) {1};
  \node[reg, anchor=north east] (l0) at (c0.south east) {--};
  \node[hd, anchor=north] (r0l) at (r0.south) {R1};
  \node[hd, anchor=north] (l0l) at (l0.south) {\iflangja{リンクレジスタ}{link register}};
  % ---- core 1
  \node[core] (c1) at (8.4,0) {\iflangja{コア}{core}~1};
  \node[reg, anchor=north west] (r1) at (c1.south west) {0};
  \node[reg, anchor=north east, fill=ln-caution!15] (l1) at (c1.south east) {--};
  \node[hd, anchor=north] (r1l) at (r1.south) {R1};
  \node[hd, anchor=north] (l1l) at (l1.south) {\iflangja{リンクレジスタ}{link register}};
  % ---- operations, numbered in the order they happen
  \node[op] at ([yshift=2pt]c0.north west) {%
    \iflangja{(1) LL: R1 $=$ 0，リンクレジスタ $\leftarrow$ A}%
             {(1) LL: R1 $=$ 0, link register $\leftarrow$ A}\\
    \iflangja{(3) SC: リンク有効 $\to$ 1 を格納，R1 $=$ 1}%
             {(3) SC: link valid $\to$ store 1, R1 $=$ 1}};
  \node[op] at ([yshift=2pt]c1.north west) {%
    \iflangja{(2) LL: R1 $=$ 0，リンクレジスタ $\leftarrow$ A}%
             {(2) LL: R1 $=$ 0, link register $\leftarrow$ A}\\
    \textcolor{ln-caution}{\iflangja{(5) SC: リンク無効 $\to$ 格納せず，R1 $=$ 0}%
             {(5) SC: link cleared $\to$ no store, R1 $=$ 0}}};
  % ---- memory
  \node[big, minimum width=4.4cm] (mem) at (5.1,-2.55)
    {\iflangja{ロック変数（アドレス A）$=$ 1}{lock variable (address A) $=$ 1}};
  % store of core 0
  \draw[->] (r0l.south) |- (mem.west)
    node[pos=0.75, above, font=\scriptsize] {\iflangja{(3) 1 を格納}{(3) store 1}};
  % clearing of the link registers
  \draw[->, ln-caution] (mem.east) -- (11.0,-2.55) |- (l1.east)
    node[pos=0, above left, font=\scriptsize, align=right, inner sep=2pt]
    {\iflangja{(4) A への書き込みで，\\A を保持する\\リンクレジスタを消去}%
              {(4) the write to A clears\\every link register holding A}};
\end{tikzpicture}

  \caption{LL/SC on two cores.  The successful store of core 0 clears the
    link register of core 1, whose store-conditional then fails.}
  \label{fig:l20-llsc}
\end{figure}


\needspace{6\baselineskip}
\begin{example}[Two cores competing for a lock]
  \label{ex:l20-llsc}
  Cores 0 and 1 execute the \texttt{lock} code above at the same time; the
  lock variable, at address A, is 0.  \Cref{tab:l20-llsc} traces the
  instructions in the order in which they execute.  Both cores find the lock
  free, but only one \texttt{SC} can succeed: its store clears the other
  core's link register.  Core 1 goes back to \texttt{LL}, now reads 1, and
  waits.
\end{example}

\begin{table}[htbp]
  \centering\small
  \caption{Trace of \cref{ex:l20-llsc}.  A dash in a link-register column
    means that the register holds no address.}
  \label{tab:l20-llsc}
  \begin{tabular}{@{}ccllccc@{}}
    \toprule
    & & & & \multicolumn{2}{c}{Link register} & Lock \\
    \cmidrule(lr){5-6}
    Step & Core & Instruction & Effect & Core 0 & Core 1 & variable \\
    \midrule
    1  & 0 & \texttt{LL R1, 0(R2)}  & R1 $=$ 0                & A  & -- & 0 \\
    2  & 1 & \texttt{LL R1, 0(R2)}  & R1 $=$ 0                & A  & A  & 0 \\
    3  & 0 & \texttt{BNEZ}, \texttt{ADDI} & R1 $=$ 1          & A  & A  & 0 \\
    4  & 1 & \texttt{BNEZ}, \texttt{ADDI} & R1 $=$ 1          & A  & A  & 0 \\
    5  & 0 & \texttt{SC R1, 0(R2)}  & stores 1, R1 $=$ 1      & -- & -- & 1 \\
    6  & 1 & \texttt{SC R1, 0(R2)}  & no store, R1 $=$ 0      & -- & -- & 1 \\
    7  & 0 & \texttt{BEQZ R1, lock} & not taken: holds lock   & -- & -- & 1 \\
    8  & 1 & \texttt{BEQZ R1, lock} & taken: start over       & -- & -- & 1 \\
    9  & 1 & \texttt{LL R1, 0(R2)}  & R1 $=$ 1                & -- & A  & 1 \\
    10 & 1 & \texttt{BNEZ R1, lock} & taken: lock is held     & -- & A  & 1 \\
    \bottomrule
  \end{tabular}
\end{table}

\begin{keypoint}
  A lock needs the read and the write of the lock variable to be atomic.
  Test-and-set does both in one unusual instruction; LL/SC uses two ordinary
  memory accesses and lets the store fail if anyone wrote the location in
  between.
\end{keypoint}

\section{Locks and performance}
\label{sec:l20-performance}

The lock variable is a shared variable like any other: every core that uses
the lock has the block containing it in its cache, and the cache coherence
protocol of Lesson 19 keeps these copies coherent.\source{Lesson 20, videos
14--17; H\&P §5.5.}  A spin lock built on test-and-set interacts badly with
coherence.

\texttt{TST} writes the lock variable even when the lock is held; writing 1
over 1 changes no value, but for the coherence protocol it is a write like
any other.  Under the MSI protocol the writing core must have the block in
the modified state, and all other copies are invalidated.  While one core
holds the lock, each waiting core that executes \texttt{TST} takes the block
away from the core that tried before it.  Every attempt of every waiting
core thus causes a coherence miss and a transfer of the block over the bus
(\cref{tab:l20-coherence}, top).  The waiting cores accomplish nothing, yet
their bus traffic slows down the memory accesses of all other cores,
including the lock holder executing its critical section, which also has to
fetch the block back when it unlocks.

\needspace{7\baselineskip}
The remedy is to test the lock variable with an ordinary load first and to
use the atomic instruction only when the lock appears free:
\begin{lstlisting}[language={[x86masm]Assembler}, numbers=none, deletekeywords={lock}]
lock:   LW   R1, 0(R2)   ; ordinary read of lock variable
        BNEZ R1, lock    ; held: keep reading
        TST  R1, 0(R2)   ; looks free: try to take it
        BNEZ R1, lock    ; another thread won: start over
\end{lstlisting}
Waiting cores now only read the lock variable, so all of them keep the block
in the shared state and their loads are cache hits: while the lock is held,
the waiting cores cause no bus traffic at all
(\cref{tab:l20-coherence}, bottom).  A write, and with it an invalidation,
happens only when the lock is released or acquired.

\begin{table}[htbp]
  \centering\small
  \caption{MSI state of the block holding the lock variable in the caches of
    four cores (M, S, I: modified, shared, invalid) while core 0 holds the
    lock.  Top: cores 1--3 spin on
    \texttt{TST}.  Bottom: they test with \texttt{LW} first.}
  \label{tab:l20-coherence}
  \begin{tabular}{@{}lccccl@{}}
    \toprule
    Access & C0 & C1 & C2 & C3 & Bus traffic \\
    \midrule
    \multicolumn{6}{@{}l}{\emph{Spinning on test-and-set}} \\
    start: C0 holds the lock                  & M & I & I & I & \\
    C1: \texttt{TST} (reads 1)                & I & M & I & I & miss, block moves \\
    C2: \texttt{TST} (reads 1)                & I & I & M & I & miss, block moves \\
    C3: \texttt{TST} (reads 1)                & I & I & I & M & miss, block moves \\
    C1: \texttt{TST} (reads 1)                & I & M & I & I & miss, block moves \\
    \quad\dots\ one miss per attempt, until   &   &   &   &   & \\
    C0: \texttt{SW R0} (unlock)               & M & I & I & I & miss, block moves \\
    \midrule
    \multicolumn{6}{@{}l}{\emph{Testing first}} \\
    start: C0 holds the lock                  & M & I & I & I & \\
    C1: \texttt{LW} (reads 1)                 & S & S & I & I & miss \\
    C2: \texttt{LW} (reads 1)                 & S & S & S & I & miss \\
    C3: \texttt{LW} (reads 1)                 & S & S & S & S & miss \\
    C1--C3: \texttt{LW}, repeatedly           & S & S & S & S & none (hits) \\
    C0: \texttt{SW R0} (unlock)               & M & I & I & I & invalidation \\
    C2: \texttt{LW} (reads 0)                 & S & I & S & I & miss \\
    C2: \texttt{TST} (reads 0, acquires)      & I & I & M & I & invalidation \\
    \bottomrule
  \end{tabular}
\end{table}

Two points about the test-first lock are easy to get wrong.
\begin{itemize}
  \item \textbf{The load is only a filter.}  Mutual exclusion still comes
    from the atomic instruction.  Between the \texttt{LW} that finds the
    lock free and the \texttt{TST}, another thread may take the lock, so the
    value returned by \texttt{TST} must still be tested.  A lock that
    branches into the critical section on the result of the \texttt{LW}
    alone is the broken lock of \cref{sec:l20-implementing}.
  \item \textbf{LL/SC already tests first.}  \texttt{LL} is a read, and the
    lock of \cref{sec:l20-llsc} attempts the \texttt{SC} only after
    \texttt{LL} has found the lock free.  A waiting core that keeps finding
    1 never writes, so no extra load is needed.
\end{itemize}

\texttt{unlock}, finally, needs no atomic instruction: an ordinary store of
0 suffices.  Only the thread that holds the lock executes \texttt{unlock},
so two threads never release the same lock at the same time.  A thread
trying to acquire the lock at that moment reads either 1, before the store,
and keeps waiting, or 0, after it, and competes for the lock---both are
correct outcomes.  With LL/SC the store also clears the link registers of
cores that have just load-linked the variable; these cores read 1 and do not
attempt an \texttt{SC}, so they are not affected.

\begin{keypoint}
  Correctness requires the write to the lock variable to be atomic;
  performance requires it to be rare.  Waiting threads should only read the
  lock variable, so that its block stays shared in all caches.
\end{keypoint}

\section{Barrier synchronization}
\label{sec:l20-barrier}

Locks make threads take turns.\source{Lesson 20, videos 18--22; H\&P
App.~I.}  A different form of coordination makes threads wait for each
other.

\begin{definition}[Barrier synchronization]
  In \term{barrier synchronization}[バリア同期] a group of $N$ threads calls
  a \emph{barrier} at a common point in the program.  A thread that reaches
  the barrier waits there; when the last of the $N$ threads arrives, all of
  them continue (\cref{fig:l20-barrier}).
\end{definition}

\begin{figure}[htbp]
  \centering
  % Barrier synchronization with three threads over two rounds.  A thread that
% reaches the barrier waits until the last thread of its round arrives; then
% all three continue together.
\begin{tikzpicture}[x=0.85cm,y=1cm, font=\footnotesize\sffamily, >={Stealth[length=1.6mm]}]
  \def\lnbox#1#2#3#4{% y, x0, x1, fill
    \draw[fill=#4] (#2,{#1-0.2}) rectangle (#3,{#1+0.2});}
  % releases of the barrier
  \foreach \x/\r in {5/{0 $\to$ 1}, 9.5/{1 $\to$ 0}} {
    \draw[dashed, ln-accent, thick] (\x,-0.55) -- (\x,2.45);
    \node[font=\scriptsize, text=ln-accent, anchor=south, align=center] at (\x,2.45)
      {\iflangja{解放}{released}\\\texttt{release}: \r};
  }
  % thread/y/arrival in round 0/arrival in round 1
  \foreach \t/\y/\a/\b in {0/2/2.5/9.5, 1/1/4/7.5, 2/0/5/8} {
    \node[anchor=east] at (-0.15,\y) {\iflangja{スレッド}{thread}~\t};
    \lnbox{\y}{0}{\a}{white}
    \ifdim\a pt<5pt \lnbox{\y}{\a}{5}{ln-caution!15}\fi
    \lnbox{\y}{5}{\b}{white}
    \ifdim\b pt<9.5pt \lnbox{\y}{\b}{9.5}{ln-caution!15}\fi
    \lnbox{\y}{9.5}{10.5}{white}
    \fill (\a,{\y+0.2}) -- ++(-0.09,0.16) -- ++(0.18,0) -- cycle;
    \fill (\b,{\y+0.2}) -- ++(-0.09,0.16) -- ++(0.18,0) -- cycle;
  }
  % time axis
  \draw[->, ln-muted] (0,-0.55) -- (10.8,-0.55)
    node[right, inner sep=2pt] {\iflangja{時間}{time}};
  % legend
  \begin{scope}[yshift=-1.0cm]
    \fill (0.09,-0.02) -- ++(-0.09,0.16) -- ++(0.18,0) -- cycle;
    \node[anchor=west, font=\scriptsize] at (0.25,0.05) {\iflangja{バリアに到着}{reaches the barrier}};
    \draw[fill=white] (3.2,-0.07) rectangle ++(0.35,0.25);
    \node[anchor=west, font=\scriptsize] at (3.6,0.05) {\iflangja{計算}{computing}};
    \draw[fill=ln-caution!15] (5.4,-0.07) rectangle ++(0.35,0.25);
    \node[anchor=west, font=\scriptsize] at (5.8,0.05) {\iflangja{バリアで待機}{waiting at the barrier}};
  \end{scope}
\end{tikzpicture}

  \caption{A barrier used in two successive rounds by three threads.  Each
    round is released when its last thread arrives.}
  \label{fig:l20-barrier}
\end{figure}

Barriers suit parallel programs that proceed in steps, where every thread
computes its part of one step and the next step needs the results of all
threads.  The shared-memory sum of Lesson 18 is an example: thread 0 may
print the total only after all threads have added their partial sums, that
is, after a barrier.

\subsection{A simple implementation}

A shared counter records how many threads have arrived, and each thread
waits until it reaches $N$:
\begin{lstlisting}[language=C, numbers=none, morekeywords={shared}]
shared int    count = 0;
shared lock_t countLock;
void barrier(void) {
  lock(&countLock);
  count++;              /* one more thread arrived */
  unlock(&countLock);
  while (count < N)
    ;                   /* wait for the others     */
}
\end{lstlisting}
The increment needs the lock: a lost increment (\cref{tab:l20-race}) would
keep \texttt{count} below $N$ forever.  The waiting loop may read
\texttt{count} without the lock, because the counter only grows; a thread
that reads a stale value merely spins once more.

\subsection{Why the simple implementation fails}

The simple barrier works once.  A barrier, however, is normally called in
every step of a loop, and in the second step \texttt{count} is already $N$:
it keeps growing past $N$, no thread waits, and the barrier has no effect.
Resetting the counter is not a simple fix.
\begin{itemize}
  \item If the thread that brings the counter to $N$ also resets it to 0,
    the condition $\texttt{count} = N$ exists only for a moment.  A waiting
    thread that happens to read the counter only after the reset sees
    $0 < N$ and keeps waiting.  It is waiting for a round to which it
    already belongs, while the other threads go on to the next barrier;
    there only $N-1$ threads can arrive, so every thread waits forever.
  \item If each thread instead decrements the counter when it leaves, the
    counter drops below $N$ as soon as the first thread leaves, so a thread
    that has not yet noticed the release again keeps waiting; moreover, a
    fast thread can reach the next barrier before a slow thread has left the
    previous one, and the two rounds are counted together.
\end{itemize}
The underlying problem is that a waiting thread must learn that \emph{its}
round has been released, no matter how late it looks.

\subsection{A reusable barrier}

A reusable barrier therefore records every release permanently.  It
numbers the rounds: under the lock, each thread notes the number of the
round it joins, and the thread that completes a round both resets the
counter and advances the round number.  Each thread then waits until the
round number differs from the one it noted:
\needspace{14\baselineskip}
\begin{lstlisting}[language=C, numbers=none, morekeywords={shared}]
shared int    count = 0, round = 0;
shared lock_t countLock;
void barrier(void) {
  lock(&countLock);
  int myRound = round;    /* the round I join     */
  if (++count == N) {     /* I am the last thread */
    count = 0;            /* start the next round */
    round++;              /* release this one     */
  }
  unlock(&countLock);
  while (round == myRound)
    ;                     /* wait for the release */
}
\end{lstlisting}

Joining a round, counting and releasing all happen under the same lock, so
they are consistent with each other.  The release changes \texttt{round}
for good: a thread that looks late still sees that its round is over, and
the last thread of a round does not wait at all.  A thread that arrives
after a release reads the new round number and is counted in the next
round, even if threads of the previous round have not yet left their
waiting loops.  As before, \texttt{round} may be read without the lock,
because it only grows and any value other than \texttt{myRound} means that
the thread's round has been released.  In \cref{fig:l20-barrier},
\texttt{round} is 0 while the first round forms, becomes 1 when thread~2
arrives, and becomes 2 when thread~0 completes the second round.

\begin{keypoint}[Lesson summary]
  Threads that update shared data must synchronize, or updates are lost
  depending on timing.  A lock provides mutual exclusion for a critical
  section, and all accesses to a variable must use the same lock.  A lock
  cannot be built from ordinary loads and stores, because its own test and
  set would need a critical section; software solutions such as Lamport's
  bakery algorithm are too slow, so processors provide atomic instructions.
  Test-and-set reads and writes memory in one instruction, which is awkward
  for a load/store pipeline; LL/SC uses two ordinary memory accesses and a
  link register that makes the store fail if the location was written in
  between.  A spin lock should test the lock variable with a plain read
  before its atomic write, so that waiting cores keep the block shared
  instead of moving it between caches; unlock is a plain store.  Barrier
  synchronization makes a group of threads wait for each other; a barrier
  that is used repeatedly must record each release permanently, for example
  as a round number, rather than rely on a counter that is reset.
\end{keypoint}

\end{document}
